%! Unit @@UNITINT@@: @@UNITTITLE@@ — Summative Assessment — @@TESTKIND@@ (blank)
\documentclass[10pt]{article}
\usepackage{@@PREFIX@@-article}
\usepackage{@@PREFIX@@-boxes}

% Part-divider strip (headlinebox takes one arg: the background color)
\newcommand{\parthead}[1]{%
  \vspace{6pt}\noindent
  \begin{headlinebox}{navy}\color{white}\bfseries #1\end{headlinebox}%
  \vspace{4pt}}

\begin{document}

% The unit test is taken in a testing setting, so it DOES carry the name row —
% it is the one exception to the namestrip rule besides the packet cover.
\pageheader{Unit @@UNITINT@@: @@UNITTITLE@@}{@@TESTKIND@@}
\namedateperiod
@@TESTINTRO@@
% ======================================================================
% Mirror the AP exam's shape: Section I multiple choice, Section II free response.
% The practice test and the actual test must be the same length as their keys —
% put every worked solution in a \begin{work} block authored identically in both.
\parthead{Section I --- Multiple Choice}
% TODO: AP-style multiple-choice items in context, five options each.

\parthead{Section II --- Free Response}
% TODO: 2--3 multi-part FRQs. Leave \vspace work room (or a \begin{work} block) after each
% part, and end at least one with an interpret/justify-in-context part.

\end{document}
